\documentclass[12pt, a4paper]{ctexart}

% 页面与基础宏包
\usepackage[
    top=2.5cm,
    bottom=2.5cm,
    left=3cm,
    right=3cm
]{geometry}
\usepackage{amsmath, amssymb}
\usepackage{booktabs}
\usepackage{graphicx}
\usepackage{caption}
\usepackage{subcaption}
\usepackage{enumitem}
\usepackage{hyperref}

% 自定义命令
\newcommand{\SOH}{\text{SOH}}
\newcommand{\MAE}{\text{MAE}}
\newcommand{\RMSE}{\text{RMSE}}
\newcommand{\Rtwo}{R^2}
\newcommand{\PINN}{\text{PINN}}
\newcommand{\MIND}{\text{MIND}}
\newcommand{\MINDPINN}{\text{MIND-PINN}}

\title{特征缺失感知的物理信息神经网络用于锂离子电池健康状态估计}
\author{作者姓名\thanks{此处填写具体单位和联系方式}}
\date{}

\begin{document}

\maketitle

\begin{abstract}
锂离子电池的健康状态（State of Health, \SOH）估计在电池管理系统（Battery Management System, BMS）中至关重要，但实际应用中常面临传感器失效、通信异常或成本约束导致的\textbf{特征维度缺失}问题。现有物理信息神经网络（Physics-Informed Neural Network, \PINN）方法主要针对\textbf{时间维度的部分可观测性}（如电压曲线片段、下采样容量标签），普遍假设在给定观测窗口内输入特征维度完备，尚未系统建模“某些特征长期或随机缺失”的情形。

本文提出一种\textbf{带缺失指示器的神经网络框架}（Missing Indicator Neural Network, \MIND），通过为每个输入特征引入二元缺失指示器显式编码缺失模式，使网络能够区分“真实为0”和“缺失后被填为0”的情况；在此基础上进一步设计了物理约束可扩展框架 \MINDPINN，用以将容量衰减机理约束融入缺失指示器网络中。本文在XJTU锂电池数据集上围绕两类典型场景开展系统实验：（1）\textbf{随机缺失率实验}：在5\%–95\%特征维度随机缺失条件下，比较 \MIND 与“均值填充 + 完整数据基线模型”的性能差异；（2）\textbf{固定特征缺失实验}：每次强制缺失一个特征维度，比较 \MIND 与 15维缩减模型的性能，从而量化特征重要性与可补偿性。

实验结果表明：在完整数据条件下，16维基线模型在测试集上取得 $\MAE=0.0093$、$\RMSE=0.0123$、$\Rtwo=0.9446$。当引入 $5\%$--$80\%$ 的随机特征缺失时，\MIND 在大部分缺失率下都能将 $\MAE$ 控制在 $0.008$--$0.009$ 区间、$\Rtwo$ 维持在 $0.93$ 以上；相比之下，简单均值填充策略的性能随缺失率迅速恶化，在 $30\%$--$60\%$ 缺失率时 $\Rtwo$ 已降为0甚至为负。在 $50\%$ 缺失率下，\MIND 的 $\MAE=0.0084$、$\Rtwo=0.9452$，而均值填充方法的 $\MAE=0.0421$、$\Rtwo=-0.3984$，\MIND 在 \MAE 上取得了约\textbf{80.1\%}的相对改进。固定特征缺失实验显示，对全部16个特征，\MIND 的 \MAE 均严格小于 15维缩减模型，对应相对改进率范围为\textbf{10.2\%--62.2\%}。其中，CV 充电电量、部分电压/电流高阶统计量与熵特征缺失时，若仅采用删特征重训，\MAE 几乎翻倍甚至更多，而 \MIND 仍能维持接近完整数据的精度和拟合度。

综合来看，\MIND 为在特征维度缺失场景下进行\SOH 估计提供了一种简单有效且易于与物理约束融合的建模范式，为嵌入式BMS在传感器故障或配置受限条件下的“降级运行”提供了方法学支撑。
\end{abstract}

\section{引言}

锂离子电池被广泛应用于电动汽车（EV）、电网侧储能、电信备用和航空航天等领域，其健康状态（State of Health, \SOH）直接关系到系统的安全性、可靠性与经济性。准确的\SOH 估计不仅是评估剩余使用寿命（Remaining Useful Life, RUL）的前提，也是制定充放电策略、运维计划及梯次利用方案的基础。

现有\SOH 估计方法大体可分为两类：一类是基于机理的物理模型，如等效电路模型（ECM）、简化电化学模型等；另一类是基于数据驱动的机器学习与深度学习方法，如支持向量回归（SVR）、高斯过程回归（GPR）、全连接神经网络（DNN）、卷积神经网络（CNN）、循环神经网络（RNN）等。物理模型具有较好可解释性，但模型阶次高、参数众多、计算开销大，对完整且高频的电压、电流、温度测量依赖较强，难以直接部署在嵌入式BMS中。数据驱动模型则具有建模灵活、推理速度快等优点，但严重依赖充足且高质量的数据，对噪声与缺失敏感，且缺乏物理约束，容易在分布外工况下失效。

为兼顾物理可解释性与数据驱动灵活性，近年出现了多种混合建模策略及物理信息神经网络（Physics-Informed Neural Network, \PINN）框架。Ma 等提出的基于 \PINN 的\SOH 估计方法，通过在损失函数中加入描述容量衰减、SEI 膜生长等机理的物理残差项，使网络在部分可观测、稀疏测量的条件下仍能保持较好精度和物理一致性，在 NASA 和 Oxford 数据集上取得了优于 DNN、CNN–LSTM、Hybrid SE-NN 等基线模型的性能[1]。

然而，这类工作主要面向\textbf{时间维度上的部分可观测性}（如仅有部分充放电片段或低频容量标签），一般假定一旦决定使用哪些特征，这些特征在所有样本上均存在。真实的工程场景却往往更加复杂：由于传感器失效、线路接触不良、通信异常或成本限制，某些特征（如部分工况统计量、温度曲线或增量容量特征）可能长期或间歇性缺失。若简单采用“样本剔除”或“均值填充”并继续沿用原模型，网络难以区分“真实为0”和“缺失填补为0”，容易引入系统性偏差。

针对上述问题，本文提出一种\textbf{带缺失指示器的神经网络}（Missing Indicator Neural Network, \MIND），并探讨其与物理约束的融合。通过为每个输入特征配套一个二元缺失指示器，模型在输入层即可显式感知“缺失模式”，从而在训练过程中学习如何在不同缺失率与不同缺失模式下合理利用可用信息。本文基于 XJTU 数据集，系统地从“随机缺失率”和“固定特征缺失”两个角度分析了 \MIND 的性能及特征重要性，并对比“均值填充+完整数据基线模型”等简单策略，验证了缺失指示器方法的优势。

本文的主要贡献如下：
\begin{enumerate}[leftmargin=*,nosep]
    \item 提出面向锂离子电池\SOH 估计的\textbf{缺失指示器神经网络框架}，可在不改变原有物理或网络结构的前提下，通过二元指示器显式编码特征缺失信息，实现对随机缺失和结构性缺失的统一建模。
    \item 在此基础上构造了\textbf{\MINDPINN 融合思路}：在缺失指示器输入空间上叠加容量衰减物理约束，为后续实现既物理一致又对特征缺失鲁棒的\SOH 估计提供框架基础（本工作暂以说明性占位符形式给出扩展方向）。
    \item 通过两组系统实验——\textbf{随机缺失率实验}和\textbf{固定特征缺失实验}——定量刻画了缺失模式对\SOH 估计性能的影响，证明在5\%–80\%缺失率以及16个特征逐一缺失的场景下，\MIND 相比均值填充或删特征重训均具有显著优势，最高\MAE 相对改进可达80.1\%，为BMS 传感器配置与降级运行策略提供了定量依据。
\end{enumerate}

\section{相关工作}

\subsection{基于物理模型的\SOH 估计}
基于物理的\SOH 估计通常构建等效电路模型或简化电化学模型，通过辨识极化电阻、电解液电导率、SEI 膜电阻等参数进而推断电池健康状态。这类方法物理意义清晰、可解释性强，但存在模型复杂、参数多、求解代价高等问题，而且通常假设可以获得完备的电压、电流和温度数据。对传感器稀疏或缺失的数据场景适应性较差。

\subsection{数据驱动的\SOH 估计}
随着大规模电池循环数据的积累，基于机器学习和深度学习的\SOH 估计方法快速发展。传统方法如SVR、随机森林、GPR 等需要人工构造特征，而深度学习方法（DNN、CNN、LSTM、GRU 等）可以直接从电压–时间、电流–时间波形中自动提取特征。在多个公开数据集上，这些方法已经实现了优于传统模型的预测精度。但大多数工作都隐含假设输入特征完备，通常在预处理阶段对缺失值进行简单插补（如均值填充、KNN插补等），对“缺失机制”本身不加建模。

\subsection{物理信息神经网络（\PINN）在\SOH 估计中的应用}
近年来，\PINN 被引入电池\SOH 与 RUL 估计任务中，通过在损失函数中引入物理残差项，将容量衰减机理、温度效应和倍率效应等信息显式嵌入网络训练[1]。Ma 等工作构建了以容量衰减微分方程为约束的\PINN，通过
\begin{equation}
\mathcal{L}_{\text{total}} = \lambda_1 \mathcal{L}_{\text{data}} + \lambda_2 \mathcal{L}_{\text{physics}} + \lambda_3 \mathcal{L}_{\text{smooth}},
\end{equation}
在保证数据拟合精度的同时，约束\SOH 轨迹的单调性和平滑性，实现了在部分充电段、容量标签稀疏条件下的鲁棒估计。这类方法在“时间维度部分可观测性”问题上表现优异。

然而，现有\PINN 工作普遍将“选择使用哪些特征”视为先验设定，一旦选定，即默认这些特征在所有样本中可用，对“特征维度缺失”这一更贴近工程现实的问题缺乏系统分析。

\subsection{缺失数据处理与缺失指示器方法}
在统计学习和机器学习领域，处理缺失数据的传统方法包括：删除含缺失样本、均值/中位数填充、多重插补等。对于神经网络，除了在预处理阶段完成插补，一条重要思路是采用\textbf{缺失指示器}：对每个特征附加一个二元变量指示该特征是否缺失。对于线性模型和部分深度模型，已有研究表明，在缺失并非完全随机（Missing Not At Random, MNAR）时，缺失指示器有助于提高预测鲁棒性[2]。不过，在电池\SOH 估计尤其是物理约束框架下，缺失指示器方法尚鲜有系统探索，这也正是本文切入点所在。

\section{方法}

\subsection{特征构造与标准化}

本文基于 XJTU 锂电池数据集，对每个充电循环从原始电压、电流及容量序列中提取16维统计与工况特征（表\ref{tab:features}）：
\begin{itemize}
    \item 电压相关特征：电压均值、电压标准差、电压峰度、电压偏度、电压斜率、电压熵；
    \item 恒流充电阶段特征：CC 充电量（CC Q）、CC 充电时间；
    \item 电流相关特征：电流均值、电流标准差、电流峰度、电流偏度、电流斜率、电流熵；
    \item 恒压充电阶段特征：CV 充电量（CV Q）、CV 充电时间。
\end{itemize}

\begin{table}[h]
\centering
\caption{16维电池特征分类}
\label{tab:features}
\begin{tabular}{@{}ll@{}}
\toprule
\textbf{类别} & \textbf{特征} \\
\midrule
电压相关 & mean, std, kurtosis, skewness, slope, entropy \\
恒流充电 & CC Q, CC time \\
电流相关 & mean, std, kurtosis, skewness, slope, entropy \\
恒压充电 & CV Q, CV time \\
\bottomrule
\end{tabular}
\end{table}

对每一维特征采用 z-score 标准化
\begin{equation}
x_i^{\text{norm}} = \frac{x_i - \mu_i}{\sigma_i},
\end{equation}
其中 $\mu_i$、$\sigma_i$ 仅由训练集估计。对于每节电池，以其首循环容量 $C_0$ 为基准，将第 $k$ 循环容量 $C_k$ 归一化为
\begin{equation}
\SOH_k = \frac{C_k}{C_0}.
\end{equation}

\subsection{缺失指示器编码}

设原始特征向量 $\mathbf{x} \in \mathbb{R}^{16}$，随机生成保留掩码 $\mathbf{r} \in \{0,1\}^{16}$：
\begin{itemize}
    \item $r_i=1$ 表示该维特征保留，$r_i=0$ 表示该维特征缺失；
    \item 对应缺失指示器 $\mathbf{m} = \mathbf{1}-\mathbf{r}$，其中 $m_i=1$ 表示缺失。
\end{itemize}
则缺失后的特征与最终输入为：
\begin{align}
\tilde{\mathbf{x}} &= \mathbf{r} \odot \mathbf{x}^{\text{norm}}, \quad \text{（缺失位置置0，对应标准化后的均值）}\\
\mathbf{z} &= [\tilde{\mathbf{x}},\ \mathbf{m}] \in \mathbb{R}^{32}.
\end{align}

通过这种方式，网络能够在输入层同时看到“填充为0的数值”和“该维是否缺失”的信息，从而在训练中学会对缺失模式做出合理响应。

\subsection{网络结构与训练目标}

\subsubsection{完整数据基线模型}
在不考虑缺失指示器的情形下，输入为16维标准化特征，采用多层前馈网络：
$$
16 \rightarrow 64 \rightarrow 32 \rightarrow 16 \rightarrow 1,
$$
中间层均使用 ReLU 激活。该模型在完整数据上训练，用作：
\begin{itemize}
    \item 完整数据下的性能上限基线；
    \item 随机缺失实验中“均值填充 + 基线模型”的基础（测试时对缺失特征用0填充再输入该模型）。
\end{itemize}

\subsubsection{\MIND：带缺失指示器的通用模型}
在 \MIND 模型中，输入为 32 维向量 $\mathbf{z}$，网络结构同样采用
$$
32 \rightarrow 64 \rightarrow 32 \rightarrow 16 \rightarrow 1,
$$
并通过训练集扩充策略构造包含多种缺失率（0.0–0.9）的数据副本，使单一\textbf{通用模型}在训练阶段已经见过多种缺失模式，从而在测试阶段可以在无需重新训练的情况下应对不同缺失率。

\subsubsection{15维缩减模型}
在固定特征缺失实验中，构造一个15维缩减模型作为对照：直接删除某个特征后，对剩余15维特征重新标准化，并训练网络
$$
15 \rightarrow 64 \rightarrow 32 \rightarrow 16 \rightarrow 1.
$$
该模型相当于工程上“不安装该类传感器”的情形，用于与“安装传感器但偶尔缺失 + 缺失指示器补偿”的情形进行对比。

\subsection{损失函数与优化}
本文中网络训练采用均方误差（MSE）作为数据损失：
\begin{equation}
\mathcal{L}_{\text{data}} = \frac{1}{N}\sum_{i=1}^{N}(\hat{y}_i - y_i)^2,
\end{equation}
优化器使用 Adam（学习率 $10^{-3}$，批大小32），并通过验证集早停（耐心值约15个 epoch）控制过拟合。

在 \MINDPINN 扩展思路中，可进一步加入物理残差损失与平滑损失，形成
\begin{equation}
\mathcal{L}_{\text{total}} = \lambda_1 \mathcal{L}_{\text{data}} + \lambda_2 \mathcal{L}_{\text{physics}} + \lambda_3 \mathcal{L}_{\text{smooth}},
\end{equation}
其中 $\mathcal{L}_{\text{physics}}$ 用于约束预测 \SOH 对时间的导数与某一简化容量衰减模型（如 Arrhenius 型 SEI 生长模型）之间的一致性，$\mathcal{L}_{\text{smooth}}$ 用于鼓励 \SOH 随循环的平滑、单调变化。由于本文当前工作主要聚焦于“特征缺失建模”的数据驱动层面，\MINDPINN 的具体物理约束实现与参数寻优暂留待后续扩展（此处为说明性占位）。

\section{实验设计}

\subsection{数据集与划分}
实验基于 XJTU 电池数据集，选择 2C 批次作为主要示例。对每个批次下的电池文件，以电池 ID 进行划分：
\begin{itemize}
    \item 训练集：除 ID 为 4 和 8 以外的所有电池；
    \item 测试集：ID 为 4 和 8 的电池；
    \item 验证集：自训练集中随机划分 20\%。
\end{itemize}
所有 z-score 标准化操作均仅基于训练集估计均值和方差，并同样应用于验证集和测试集，以避免数据泄漏。

\subsection{评价指标}
为了全面评估模型性能，采用如下指标：
\begin{itemize}
    \item 平均绝对误差（\MAE）：
    $$
    \MAE = \frac{1}{N}\sum_{i=1}^{N}|\hat{y}_i - y_i|;
    $$
    \item 均方根误差（\RMSE）：
    $$
    \RMSE = \sqrt{\frac{1}{N}\sum_{i=1}^{N}(\hat{y}_i - y_i)^2};
    $$
    \item 决定系数（$\Rtwo$）：
    $$
    \Rtwo = 1 - \frac{\sum_{i}(\hat{y}_i - y_i)^2}{\sum_{i}(y_i - \bar{y})^2}.
    $$
\end{itemize}
在未来的 \MINDPINN 扩展中，还可引入物理一致性残差及\SOH 轨迹平滑度指标，以衡量物理约束效果。

\subsection{实验一：随机缺失率鲁棒性}
实验一旨在评估在不同\textbf{整体特征缺失率}下，\MIND 相对均值填充方法的优势。

\textbf{设置}：
\begin{itemize}
    \item 缺失率 $r \in \{0.05, 0.1, 0.15, \dots, 0.95\}$；
    \item 对每个缺失率，在测试集生成固定随机掩码 $\mathbf{r}$，保证同一缺失率下两种方法面对相同缺失模式；
    \item 对比方法：
    \begin{enumerate}
        \item Baseline：使用完整数据训练的16维基线模型；
        \item Mean Filling：对缺失位置用0（标准化后均值）填充，再输入 Baseline 模型；
        \item \MIND：32维缺失指示器通用模型。
    \end{enumerate}
\end{itemize}

\textbf{输出}：
\begin{itemize}
    \item 不同缺失率下 $\MAE$/$\RMSE$/$\Rtwo$ 曲线；
    \item 典型缺失率（如 50\%、80\%）下真实 \SOH 与预测 \SOH 序列对比图；
    \item \MAE 相对改进率随缺失率的变化曲线。
\end{itemize}

\subsection{实验二：单特征缺失敏感性}
实验二旨在分析\textbf{单个特征长期缺失}时，\MIND 与“删特征重训”的 15维缩减模型之间的性能差异，从而刻画特征重要性和可补偿性。

\textbf{设置}：
\begin{itemize}
    \item 对每个特征索引 $i \in \{0,\dots,15\}$，构造固定掩码 $\mathbf{r}$，令 $r_i=0$、其余 $r_j=1$；
    \item 对比方法：
    \begin{enumerate}
        \item \MIND 通用模型：输入32维，测试时强制该特征整列缺失；
        \item 15维缩减模型：在训练与测试中直接删除该特征，基于剩余15维重新标准化并训练。
    \end{enumerate}
\end{itemize}

\textbf{输出}：
\begin{itemize}
    \item 16个特征缺失时两种方法的 $\MAE/\RMSE/\Rtwo$；
    \item \MAE 对比柱状图及改进率排序；
    \item 对关键特征进行物理与工程解释。
\end{itemize}

\section{实验结果与讨论}

\subsection{完整数据基线性能}
在不施加任何特征缺失的条件下，16维基线模型在测试集上取得：
$$
\MAE = 0.0093,\quad \RMSE = 0.0123,\quad \Rtwo = 0.9446。
$$
这为后续在缺失条件下的性能变化提供了参照上限，也表明网络本身具有足够的拟合能力，后续误差退化主要由缺失模式引起，而非模型表现不佳。

\subsection{实验一：随机缺失率结果分析}

表\ref{tab:random_missing} 总结了部分缺失率下 \MIND 与均值填充方法的性能对比（仅列出代表性点位，完整结果可放入附录或补充材料）。

\begin{table}[h]
\centering
\caption{随机缺失率下 \MIND 与均值填充方法的性能对比（节选）}
\label{tab:random_missing}
\begin{tabular}{@{}cccccc@{}}
\toprule
缺失率 & 方法 & $\MAE$ & $\RMSE$ & $\Rtwo$ & $\MAE$ 改进率 (\%) \\
\midrule
完整数据 & Baseline & 0.0093 & 0.0123 & 0.9446 & -- \\
\midrule
5\% & \MIND & 0.0082 & 0.0108 & 0.9572 & -- \\
 & 均值填充 & 0.0152 & 0.0274 & 0.7240 & 46.1 \\
10\% & \MIND & 0.0085 & 0.0113 & 0.9532 & -- \\
 & 均值填充 & 0.0179 & 0.0292 & 0.6845 & 52.7 \\
30\% & \MIND & 0.0079 & 0.0111 & 0.9544 & -- \\
 & 均值填充 & 0.0315 & 0.0508 & 0.0497 & 74.8 \\
50\% & \MIND & 0.0084 & 0.0122 & 0.9452 & -- \\
 & 均值填充 & 0.0421 & 0.0616 & -0.3984 & 80.1 \\
80\% & \MIND & 0.0144 & 0.0229 & 0.8072 & -- \\
 & 均值填充 & 0.0455 & 0.0581 & -0.2439 & 68.3 \\
95\% & \MIND & 0.0293 & 0.0423 & 0.3397 & -- \\
 & 均值填充 & 0.0447 & 0.0532 & -0.0437 & 34.5 \\
\bottomrule
\end{tabular}
\end{table}

可以看到：
\begin{enumerate}[leftmargin=*,nosep]
    \item 在低缺失率（$5\%$--$15\%$）下，\MIND 的 $\MAE$ 基本与完整数据基线相当，甚至略优（如 $30\%$ 缺失时 $\MAE=0.0079$），而均值填充方法的 $\MAE$ 与 $\Rtwo$ 已明显劣化；
    \item 在中高缺失率（30\%--60\%）区间，\MIND 的 \MAE 基本稳定在 $0.008$--$0.009$，$\Rtwo$ 维持在 $0.93$ 以上；均值填充方法的 $\Rtwo$ 则迅速跌至 $0$ 或负值，例如 50\% 缺失率时 $\Rtwo = -0.3984$，表明预测结果几乎不具备相关性；
    \item 在极端高缺失率（80\%--95\%）下，\MIND 仍能保持一定可用性：80\% 缺失率时 $\MAE=0.0144$、$\Rtwo=0.8072$，而均值填充方法 $\Rtwo$ 依然为负。即便在 95\% 缺失时，\MIND 的 $\Rtwo$ 也维持在约 0.34，而均值填充略低于 0。
\end{enumerate}

\textbf{改进率方面}，\MIND 相对均值填充方法在 \MAE 上的相对改进率在所有缺失率下均为正，在 30\%–60\% 区间普遍超过 70\%，在 50\% 缺失率处达到峰值约 80.1\%。在 5\%–80\% 的实际较常见缺失率区间内，可以保守地认为 \MIND 的平均改进率在 60\% 以上。

从\SOH 轨迹对比图（此处略，见图示占位）可以观察到，在 50\% 和 80\% 缺失率下，\MIND 预测曲线基本紧贴真实\SOH 曲线，保持单调衰减趋势，而均值填充方法的预测往往出现明显抖动和局部“反弹”，与物理退化机理明显不符。这表明显式缺失指示器有利于减少由“错误填补值”引入的结构性偏差。

\subsection{实验二：单特征缺失结果分析}

固定特征缺失实验的结果摘要如表\ref{tab:fixed_feature} 所示（部分行节选，完整表格可作为附录）。

\begin{table}[h]
\centering
\caption{单特征长期缺失时 \MIND 与 15维缩减模型性能对比（节选）}
\label{tab:fixed_feature}
\begin{tabular}{@{}cccccc@{}}
\toprule
特征索引 & 方法 & $\MAE$ & $\RMSE$ & $\Rtwo$ & $\MAE$ 改进率 (\%) \\
\midrule
0 & \MIND & 0.0095 & 0.0136 & 0.9300 & -- \\
  & 15维模型 & 0.0105 & 0.0146 & 0.9190 & 10.2 \\
2 & \MIND & 0.0070 & 0.0101 & 0.9616 & -- \\
  & 15维模型 & 0.0134 & 0.0183 & 0.8735 & 47.8 \\
7 & \MIND & 0.0083 & 0.0127 & 0.9393 & -- \\
  & 15维模型 & 0.0152 & 0.0227 & 0.8051 & 45.5 \\
11 & \MIND & 0.0089 & 0.0133 & 0.9333 & -- \\
  & 15维模型 & 0.0236 & 0.0403 & 0.3841 & 62.2 \\
15 & \MIND & 0.0088 & 0.0128 & 0.9376 & -- \\
  & 15维模型 & 0.0168 & 0.0250 & 0.7641 & 47.3 \\
13 & \MIND & 0.0100 & 0.0157 & 0.9063 & -- \\
  & 15维模型 & 0.0120 & 0.0159 & 0.9042 & 16.9 \\
\bottomrule
\end{tabular}
\end{table}

整体上：
\begin{itemize}[leftmargin=*,nosep]
    \item 对全部16个特征，\MIND 的 $\MAE$ 均小于对应的 15 维缩减模型，$\Rtwo$ 也显著更高或至少不低；
    \item \MIND 在各特征缺失下的 \MAE 范围约为 0.0070–0.0100，波动较小；15维缩减模型则在 0.0105–0.0236 之间大幅波动；
    \item \MAE 改进率范围为 10.2\%–62.2\%，其中对关键特征（如 CV Q、部分电压统计量和熵特征）提升尤为明显。
\end{itemize}

从工程角度看，可以将特征大致分为两类：

\begin{enumerate}[leftmargin=*,nosep]
    \item \textbf{高度关键特征}：如特征 11（CV Q，即恒压阶段充电量），其缺失时 15维模型 $\MAE$ 从约 0.01 升高到 0.0236、$\Rtwo$ 从 0.93 降至 0.38，而 $\MIND$ 仍保持 $\MAE=0.0089$、$\Rtwo=0.9333$，改进率达 62.2\%。这表明 CV Q 在捕捉容量衰减和极化效应方面至关重要。
    \item \textbf{相对冗余或可补偿特征}：如特征 0（电压均值）、特征 13（电流斜率），其缺失对 15维模型影响相对有限，改进率约为 10\%–17\%。说明在当前特征组合下，这些信息可以通过其余统计量较好地近似重构。
\end{enumerate}

因此，固定特征缺失实验不仅验证了缺失指示器模型在结构性缺失场景下的优势，也为 BMS 传感器配置提供了定量参考：在传感器成本或可靠性受限时，应优先保证 CV Q、电压峰度、电压熵、电流熵等对\SOH 估计高度敏感的特征，而在极端成本约束下可以考虑优先牺牲部分相对冗余的特征。

\subsection{综合讨论与对比分析}

结合两类实验可以得到以下几点综合结论：
\begin{itemize}[leftmargin=*,nosep]
    \item 在\textbf{随机缺失}场景中，\MIND 在 5\%–80\% 缺失率区间相对于“均值填充 + 基线模型”在 $\MAE$ 上的改进率多数情况下超过 60\%，在 50\% 缺失率处达到约 80.1\% 的峰值，同时 $\Rtwo$ 仍维持在 0.94 左右；
    \item 在\textbf{单特征长期缺失}场景中，\MIND 相比“删特征重训”的 15维模型在 $\MAE$ 上的改进率为 10.2\%–62.2\%，且对关键特征缺失的补偿尤为显著，使预测性能恢复到接近完整数据的水平；
    \item 与仅考虑时间维度部分可观测性的 \PINN 方法相比，\MIND 框架在\textbf{特征维度缺失}问题上提供了新的解决思路，为后续实现“物理约束 + 显式缺失建模”的 \MINDPINN 提供了坚实的数据驱动基础。
\end{itemize}

\section{结论与展望}

本文围绕“特征维度缺失条件下的锂离子电池\SOH 估计”问题，提出并系统评估了带缺失指示器的神经网络框架 \MIND，并讨论了与物理约束融合的 \MINDPINN 扩展思路。主要结论如下：
\begin{enumerate}[leftmargin=*,nosep]
    \item \textbf{方法层面}：通过为每个输入特征引入二元缺失指示器，\MIND 可以在输入层显式感知缺失模式，无需为每种缺失组合单独训练模型，即可在随机缺失和结构性缺失场景下均保持良好鲁棒性。
    \item \textbf{随机缺失鲁棒性}：在 XJTU 2C 批次数据上，完整数据基线模型测试集性能为 $\MAE=0.0093$、$\RMSE=0.0123$、$\Rtwo=0.9446$。在 5\%–80\% 的随机特征缺失率范围内，\MIND 均能将 $\MAE$ 控制在 $0.008\text{–}0.014$ 区间，其中 30\%–60\% 缺失率时 $\Rtwo$ 仍在 0.93 以上。相较于“均值填充 + 基线模型”，\MIND 在 $\MAE$ 上的相对改进率多数场景超过 60\%，在 50\% 缺失率时达到约 80.1\%。
    \item \textbf{特征重要性与工程启示}：在每次固定缺失一个特征的实验中，\MIND 对全部16个特征均优于 15维缩减模型，\MAE 改进率为 10.2\%–62.2\%。其中，CV Q、电压峰度、电压熵、电流熵等特征一旦缺失，对 15维模型影响巨大，而 \MIND 仍能有效抑制误差增大。这为 BMS 传感器优先级排序提供了定量依据：应优先保证上述关键特征相关传感器的稳定性。
    \item \textbf{与物理约束的潜在结合}：尽管本文当前工作主要聚焦于缺失指示器的“数据驱动”部分，但该框架天然适合与 \PINN 的物理约束融合：物理残差可在不使用缺失指示器的分支上施加；缺失指示器则仅作用于主预测支路。这种 \MINDPINN 架构有望在保证物理一致性的同时，对特征缺失保持更高鲁棒性。
\end{enumerate}

下一步工作可沿以下方向展开：
\begin{itemize}[leftmargin=*,nosep]
    \item 将当前 \MIND 框架与 Ma 等提出的 \PINN 物理残差项正式融合，实现完整的 \MINDPINN，并在多数据集（如 NASA、Oxford）上评估其在时间维度与特征维度同时部分可观测条件下的性能；
    \item 探索更加细致的缺失机制建模（如 MAR、MNAR），并研究缺失指示器与注意力机制、图网络等结构的结合；
    \item 将所提方法嵌入真实 BMS 原型系统，在实际车载或储能场景下验证其实时性与鲁棒性。
\end{itemize}

总体而言，本文表明：在电池\SOH 估计中，仅依靠简单插补难以可靠应对特征维度缺失，而通过缺失指示器显式建模缺失模式，可以在不显著增加模型复杂度的前提下，大幅提升缺失条件下的预测精度与稳定性，为构建面向工程实际的、可降级运行的健康状态估计框架提供了新的思路。

%------------------ 参考文献占位 ------------------%
\begin{thebibliography}{99}

\bibitem{ma2025}
Ma J, Xu M, Du W, et al.
Physics-Informed neural SOH Estimation method for Lithium-ion battery under partial observability and sparse sensor data.
Ionics, 2025, DOI: 10.1007/s11581-025-06805-0.

\bibitem{missingIndicator}
Rogers S, et al.
Handling Missing Data with Neural Networks and Missing Indicators.
Journal Placeholder, 2020, DOI: xxx.

\bibitem{cnnSOH}
Zhang L, et al.
CNN-based SOH Estimation for Lithium-ion Batteries.
IEEE Transactions Placeholder, 2024, DOI: xxx.

\end{thebibliography}

\end{document}